% GENERATED FILE. Edit canonical lessons, then run npm run generate:books.
% canonical-source-hash: fnv1a64:ac7875c8abceb0a4
% canonical-lessons: ZH-C456-zousan, ZH-C456-qiuzhu, ZH-C456-luren, ZH-C456-shizhu, ZH-C456-renling

\chapter{Lost and Found}
\label{ch:zh-lost-and-found}
\hlchaptermodality{456}

\begin{chapteropening}
\textbf{By the end of this chapter:} \emph{I can say we got separated, ask for help, ask a passer-by, and talk about the owner of a lost wallet claiming it.}

Complete the last lesson of chapter 456: I can say we got separated, ask for help, ask a passer-by, and talk about the owner of a lost wallet claiming it.
\end{chapteropening}

\section[zǒusàn]{zǒusàn — to get separated (in a crowd)}

\label{lesson:ZH-C456-zousan}



\begin{quote}
Before the new one: say the Mandarin for still valid, then the Mandarin for to apply for.
\end{quote}

\subsection*{You'll want to know: zǒusàn}
\emph{zǒusàn} — \textquotedblleft{}to get separated (in a crowd)\textquotedblright{}.

\begin{quote}
\emph{wǒ hé péngyou zài huǒchēzhàn zǒusàn le} — My friend and I lost each other at the station.
\end{quote}

Its characters come later; for now it is a word for the ear.

\begin{cousinweb}
\emph{zǒusàn} is built from \zh{走} 'to walk' and \emph{sàn} 'to scatter', both inherited from Middle Chinese, \emph{tsuwX} and \emph{sanH} (Guangyun, in Baxter's transcription), which glosses the first 'to run' and the second 'to break apart'.
\end{cousinweb}

\begin{grammarlens}[title={What you've built}]
One more word for getting around.
\end{grammarlens}

\subsection*{Your turn}
\begin{itemize}
\raggedright
  \item \emph{Say it:} \emph{zǒusàn}
  \item \emph{Say it:} \emph{zǒusàn}, once more
  \item \emph{Say it:} \emph{wǒ hé péngyou zài huǒchēzhàn zǒusàn le} — My friend and I lost each other at the station.
  \item \emph{Recall:} say the Mandarin for still valid, then the Mandarin for to apply for, then say \emph{zǒusàn} again
\end{itemize}

\begin{tcolorbox}[breakable,colback=teal!4,colframe=teal!35!black,title={Before you move on}]
What does zǒusàn mean? (\textquotedblleft{}To get separated (in a crowd)\textquotedblright{}.) Say it once more.
\end{tcolorbox}
\section[qiúzhù]{qiúzhù — to ask for help; to turn to someone for help}

\label{lesson:ZH-C456-qiuzhu}



\begin{quote}
Before the new one: say the Mandarin for to apply for, then the Mandarin for to get separated.
\end{quote}

\subsection*{You'll want to know: qiúzhù}
\emph{qiúzhù} — \textquotedblleft{}to ask for help; to turn to someone for help\textquotedblright{}.

\begin{quote}
\emph{wǒ zhǎo bu dào péngyou qù fúwùtái qiúzhù} — I can't find my friend, so I'm going to the information desk for help.
\end{quote}

Its characters come later; for now it is a word for the ear.

\begin{cousinweb}
\emph{qiúzhù} is built from \emph{qiú} 'to seek' and \emph{zhù} 'to help', both inherited from Middle Chinese, \emph{gjuw} and \emph{dzrjoH} (Guangyun, in Baxter's transcription), which glosses the first 'to search for' and the second 'to assist'.
\end{cousinweb}

\begin{grammarlens}[title={What you've built}]
One more word for getting around.
\end{grammarlens}

\subsection*{Your turn}
\begin{itemize}
\raggedright
  \item \emph{Say it:} \emph{qiúzhù}
  \item \emph{Say it:} \emph{qiúzhù}, once more
  \item \emph{Say it:} \emph{wǒ zhǎo bu dào péngyou qù fúwùtái qiúzhù} — I can't find my friend, so I'm going to the information desk for help.
  \item \emph{Recall:} say the Mandarin for to apply for, then the Mandarin for to get separated, then say \emph{qiúzhù} again
\end{itemize}

\begin{tcolorbox}[breakable,colback=teal!4,colframe=teal!35!black,title={Before you move on}]
What does qiúzhù mean? (\textquotedblleft{}To ask for help; to turn to someone for help\textquotedblright{}.) Say it once more.
\end{tcolorbox}
\section[\texorpdfstring{\zh{路人} (lùrén)}{路人 (lùrén)}]{\zh{路人} (lùrén) — passer-by; stranger in the street}

\label{lesson:ZH-C456-luren}



\begin{quote}
Before the new one: say the Mandarin for to get separated, then the Mandarin for to turn to someone for help.
\end{quote}

\subsection*{You'll want to know: \zh{路人}}
\textbf{\zh{路人}} — \emph{lùrén} — \textquotedblleft{}passer-by; stranger in the street\textquotedblright{}.

\begin{quote}
\textbf{\zh{我的地图不见了我问了一个路人}} \emph{wǒ de dìtú bú jiàn le wǒ wèn le yí ge lùrén} — My map had gone, so I asked a passer-by.
\end{quote}

\begin{cousinweb}
\textbf{\zh{路人}} (\emph{lùrén}) is built from \zh{路} 'a road' and \zh{人} 'a person', both inherited from Middle Chinese, \emph{luH} and \emph{nyin} (Guangyun, in Baxter's transcription).
\end{cousinweb}

\begin{grammarlens}[title={What you've built}]
One more word for getting around.
\end{grammarlens}

\subsection*{Your turn}
\begin{itemize}
\raggedright
  \item \emph{Say it:} \emph{lùrén}
  \item \emph{Say it:} \emph{lùrén}, once more
  \item \emph{Say it:} \emph{wǒ de dìtú bú jiàn le wǒ wèn le yí ge lùrén} — My map had gone, so I asked a passer-by.
  \item \emph{Recall:} say the Mandarin for to get separated, then the Mandarin for to turn to someone for help, then say \emph{lùrén} again
\end{itemize}

\begin{tcolorbox}[breakable,colback=teal!4,colframe=teal!35!black,title={Before you move on}]
What does \zh{路人} mean? (\textquotedblleft{}Passer-by; stranger in the street\textquotedblright{}.) Say it once more.
\end{tcolorbox}
\section[shīzhǔ]{shīzhǔ — the owner of a lost item}

\label{lesson:ZH-C456-shizhu}



\begin{quote}
Before the new one: say the Mandarin for to turn to someone for help, then the Mandarin for a passer-by.
\end{quote}

\subsection*{You'll want to know: shīzhǔ}
\emph{shīzhǔ} — \textquotedblleft{}the owner of a lost item\textquotedblright{}.

\begin{quote}
\emph{wǒmen zhǎodào le yí ge qiánbāo zhèngzài zhǎo shīzhǔ} — We found a wallet and we're looking for its owner.
\end{quote}

Its characters come later; for now it is a word for the ear.

\begin{cousinweb}
\emph{shīzhǔ} is built from \emph{shī} 'to lose' and \emph{zhǔ} 'an owner', both inherited from Middle Chinese, \emph{syit} and \emph{tsyuX} (Guangyun, in Baxter's transcription), which glosses the first 'to err' and the second 'to be in charge'.
\end{cousinweb}

\begin{grammarlens}[title={What you've built}]
One more word for getting around.
\end{grammarlens}

\subsection*{Your turn}
\begin{itemize}
\raggedright
  \item \emph{Say it:} \emph{shīzhǔ}
  \item \emph{Say it:} \emph{shīzhǔ}, once more
  \item \emph{Say it:} \emph{wǒmen zhǎodào le yí ge qiánbāo zhèngzài zhǎo shīzhǔ} — We found a wallet and we're looking for its owner.
  \item \emph{Recall:} say the Mandarin for to turn to someone for help, then the Mandarin for a passer-by, then say \emph{shīzhǔ} again
\end{itemize}

\begin{tcolorbox}[breakable,colback=teal!4,colframe=teal!35!black,title={Before you move on}]
What does shīzhǔ mean? (\textquotedblleft{}The owner of a lost item\textquotedblright{}.) Say it once more.
\end{tcolorbox}
\section[rènlǐng]{rènlǐng — to claim lost property}

\label{lesson:ZH-C456-renling}



\begin{quote}
Before the new one: say the Mandarin for a passer-by, then the Mandarin for the owner of a lost item.
\end{quote}

\subsection*{You'll want to know: rènlǐng}
\emph{rènlǐng} — \textquotedblleft{}to claim lost property\textquotedblright{}.

\begin{quote}
\emph{shīzhǔ kěyǐ qù fúwùtái rènlǐng qiánbāo} — The owner can claim the wallet at the information desk.
\end{quote}

Its characters come later; for now it is a word for the ear.

\begin{cousinweb}
\emph{rènlǐng} is built from \emph{rèn} 'to recognise' and \emph{lǐng} 'to receive', both inherited from Middle Chinese, \emph{nyinH} and \emph{ljengX} (Guangyun, in Baxter's transcription), which glosses the first 'to know' and the second 'to take charge of'; the sense 'to receive' is not given there.
\end{cousinweb}

\begin{grammarlens}[title={What you've built}]
One more word for getting around.
\end{grammarlens}

\subsection*{Your turn}
\begin{itemize}
\raggedright
  \item \emph{Say it:} \emph{rènlǐng}
  \item \emph{Say it:} \emph{rènlǐng}, once more
  \item \emph{Say it:} \emph{shīzhǔ kěyǐ qù fúwùtái rènlǐng qiánbāo} — The owner can claim the wallet at the information desk.
  \item \emph{Recall:} say the Mandarin for a passer-by, then the Mandarin for the owner of a lost item, then say \emph{rènlǐng} again
\end{itemize}

\begin{tcolorbox}[breakable,colback=teal!4,colframe=teal!35!black,title={Before you move on}]
What does rènlǐng mean? (\textquotedblleft{}To claim lost property\textquotedblright{}.) Say it once more.
\end{tcolorbox}
